\documentclass{article}
\usepackage[T1]{fontenc}
\usepackage{amsmath}
\usepackage{graphicx}
\usepackage{array}
\usepackage{booktabs}

\begin{document}

% =========================================================================
% SECTION 1: Description of the Proposed Solution
% =========================================================================

\section{Description of the Proposed Solution}
Clinical codification represents a critical administrative and clinical bottleneck in healthcare systems, requiring the translation of free-text diagnostic and procedural notes into standardized classification codes. In this work, we address the single-label multiclass clinical codification task, which maps raw, brief Spanish clinical report literals into their corresponding first-character ICD-10 categories, spanning 36 distinct classes ($0\text{--}9$ and $\text{A}\text{--}\text{Z}$). 

Our final engineering pipeline is designed as a highly specialized, sequential deep learning pipeline optimized to handle severe class imbalance, Catalan-Spanish bilingual expressions, medical jargon, and hardware VRAM limitations. The core components of our proposed architecture and text-processing workflow are detailed below.

\subsection{Sequential Pipeline Architecture}
The processing of a raw clinical literal flows through a multi-stage sequential architecture:
\begin{enumerate}
    \item \textbf{Text Preprocessing:} The raw clinical report is cleaned using a lightweight preprocessing filter implemented in the \texttt{preprocess\_text} function within our pipeline. This step casts the inputs to strings, strips leading and trailing whitespaces, and collapses multiple consecutive whitespaces into a single space:
    \begin{align}
        x_{\text{clean}} = \text{RegExSub}(\text{Strip}(\text{String}(x_{\text{raw}})), \text{`}\backslash\text{s+}\text{'}, \text{` }\text{'})
    \end{align}
    Crucially, we do not perform lowercase normalization, accent removal, or punctuation stripping. Preserving original casing and accents is essential, as the downstream specialized clinical tokenizer contains unique subword representations for these clinical forms.
    
    \item \textbf{Subword Tokenization:} The cleaned text is tokenized using the pre-trained Spanish biomedical tokenizer associated with the '\textit{PlanTL-GOB-ES/roberta-base-biomedical-clinical-es}' backbone. The sequences are padded or truncated to a highly optimized maximum length of $64$ tokens. While standard pipelines use a max length of $128$, our analysis of the dataset made us discover that the longest clinical literal in \texttt{codification\_data.csv} contains $49$ subwords. Setting the limit to $64$ reduces RAM consumption  and accelerates training speeds with zero loss of clinical information.
    
    \item \textbf{Biomedical Transformer Encoder:} The tokenized IDs are passed through the $12$-layer Spanish biomedical-clinical RoBERTa model. Standard multilingual models are trained on general web corpora and struggle to encode Spanish medical abbreviations (e.g. \texttt{VHC}, \texttt{HPB}, \texttt{DG}) and diagnostic idioms. This backbone, pre-trained on over 1B tokens of Spanish clinical notes, and medical publications, provides robust pre-built clinical contextual embeddings of dimension $d_{\text{odel}} =768$.
    
    \item \textbf{Custom Attention Pooling:} To aggregate the sequence hidden states into a single sequence vector, we replace standard $\langle\text{CLS}\rangle$ and Mean pooling with a parameterized \textbf{Attention Pooling Layer}. Standard $\langle\text{CLS}\rangle$ pooling suffers because the sequence-start token in RoBERTa is trained on general masked language modeling and fails to focus on short, highly dense medical nouns. Mean pooling assigns equal weight to non-informative terms (e.g. Catalan/Spanish prepositions like \texttt{de}, \texttt{en}). 
    
    \item \textbf{Multi-Sample Dropout Regularization:} To prevent overfitting on small clinical target sets, the pooled representation $\mathbf{v}$ is processed by a Multi-Sample Dropout head. We initialize five parallel dropout paths with a dropout probability of $p=0.15$. During training, the representation is cloned five times, passed through independent dropout masks, and fed into the shared linear classification head.
    
    Averaging the logits acts as a robust ensemble regularizer, smoothing the decision boundaries and preventing the classifier head from co-adapting to specific high-frequency clinical words.
    
    \item \textbf{Linear Classification Head:} A linear projection layer maps the averaged $768$-dimensional vector to the $36$ logits mentioned, corresponding to the ICD-10 clinical categories.
\end{enumerate}
\vspace{\baselineskip}

\subsection{Justification of the Deep Learning approach}
Our group transition from traditional statistical methods to a deep clinical transformer architecture was due to to the challenges inherent to medical narratives:
\begin{itemize}
    \item \textbf{Failure of Traditional Term Frequency Models:} Traditional statistical baselines rely on exact term matching. Clinical reports frequently use non-standard abbreviations, misspellings ,multi-word clinical concepts... For example, TF-IDF is incapable of mapping the abbreviation \texttt{HTA} to \textit{Hipertensión Arterial} or \texttt{VHC} to \textit{Hepatitis C}, treating them as entirely different features.
    
    \item \textbf{Bilingual Medical Terminology:} The dataset is characterized by a fluid, real-world mix of Spanish and Catalan terminology. A deep learning model trained on localized clinical text successfully maps these distinct lexical words to the same semantic space, but statistical methods would have required manually bilingual synonym dictionaries.
    
\vspace{\baselineskip}
\vspace{\baselineskip}
\vspace{\baselineskip}
\vspace{\baselineskip}





\end{itemize}

% =========================================================================
% SECTION 2: Experiments and Evaluation Setup
% =========================================================================
\section{Experiments and Evaluation Setup}
To assess our models and prevent overestimating performance, we established a robust validation partition and systematic training setup that clones rea clinical deployment constraints.
\vspace{\baselineskip}

\subsection{Dataset Partitioning and the "Split-Before-Augment" Protocol}
The primary clinical dataset (\texttt{codification\_data.csv}) consists of $13,702$ short clinical report literals. The classification target spans $36$ classes. Ten classes correspond to the \textbf{ICD-10 Procedure Coding System}, where digit \texttt{0} indicates surgical procedures and digit \texttt{3} indicates administrative procedures. The remaining $26$ classes  correspond to the \textbf{ICD-10 Clinical Modification (ICD-10-CM)} diagnostic categories.
We implemented a **stratified $80/20$ train/validation split** using the scikit-learn \texttt{train\_test\_split} utility, preserving class proportions across both sets:
\begin{itemize}
    \item \textbf{Stratified Training Split:} $10,961$ samples.
    \item \textbf{Clean Validation Split:} $2,741$ samples.
\end{itemize}


\subsection{Hyperparameter Configuration}
For both the baseline and advanced models, we systematically tracked and tuned several hyperparameters. Table~\ref{tab:hyperparameters} outlines the exact experimental configurations used during training.

\subsection{Evaluation Metrics and the Insufficiency of Accuracy}
Model evaluation was performed using Loss curves, Precision, Recall, and the Macro F1-Score. In our clinical codification task, Accuracy alone is a highly misleading and insufficient metric. 
\vspace{\baselineskip}

Because the primary hospital data is  imbalanced by some factors, a naive model that predicts only these three categories would mathematically achieve over $60\%$ validation accuracy while failing  to classify the other $33$ clinical classes. In a clinical setting, errors like those would have catastrophic implications for patients.

Therefore, we used the **Macro-averaged F1-Score** as our primary optimization and early stopping target:

\begin{table}[h]
\centering
\caption{Hyperparameter Configurations for Baseline and Advanced Models}
\label{tab:hyperparameters}
\begin{tabular}{lcc}
\toprule
\textbf{Hyperparameter} & \textbf{Baseline Model} & \textbf{Advanced Optimized Model} \\
\midrule
Pre-trained Backbone & \texttt{roberta-base-biomedical-clinical-es} & \texttt{roberta-base-biomedical-clinical-es} \\
Sequence Max Length ($L$) & $128$ & $64$ \\
Optimizer & AdamW & AdamW \\
Learning Rate (Backbone) & $2.0 \times 10^{-5}$ & $2.0 \times 10^{-5}$ \\
Learning Rate (Heads) & $2.0 \times 10^{-5}$ & $2.0 \times 10^{-5}$ \\
Weight Decay & $0.01$ & $0.01$ \\
Warmup Ratio & $10\%$ & $10\%$ \\
Physical Batch Size & $128$ & $32$ \\
Gradient Accumulation Steps & None & $4$ \\
Effective Batch Size & $128$ & $128$ \\
Max Training Epochs & $50$ & $50$ \\
Early Stopping Metric & Validation Accuracy & Validation Macro F1-Score \\
Early Stopping Patience & $10$ epochs & $10$ epochs \\
Encoder Frozen Layers & None (All $12$ active) & Embeddings + Bottom $6$ Layers \\
Dropout Probability ($p$) & $0.10$ & $0.15$ (5-sample Multi-Dropout) \\
Loss Criterion & Standard Cross-Entropy & Balanced Class-Weighted Cross-Entropy \\
\bottomrule
\end{tabular}
\end{table}

\vspace{\baselineskip}
\vspace{\baselineskip}

\section{Critical Analysis of Results}
Following the  implementation of our deep learning models, we made a comparative analysis of the validation and test results. Table~\ref{tab:results} breakdown of model performance.


\begin{table}[h]
\centering
\caption{Model Performance Comparison on Clinical Codification (Stratified Split)}
\label{tab:results}
\begin{tabular}{lcccc}
\toprule
\textbf{Pipeline Architecture} & \textbf{Val. Loss} & \textbf{Val. Acc.} & \textbf{Val. Macro F1} & \textbf{Kaggle Test Acc.} \\
\midrule
Baseline RoBERTa + $\langle\text{CLS}\rangle$ Pooling & $1.724$ & $0.565$ & $0.245$ & $0.558$ \\
Baseline RoBERTa + Mean Pooling & $1.682$ & $0.570$ & $0.258$ & $0.562$ \\
RoBERTa + Attention Pool + Multi-Dropout (No Aug.) & $1.104$ & $0.685$ & $0.452$ & $0.678$ \\
\textbf{Our Final Optimized Pipeline (Aug. + Weights + Freeze)} & \textbf{0.682} & \textbf{0.804} & \textbf{0.718} & \textbf{0.792} \\
\bottomrule
\end{tabular}
\end{table}


\subsection{Successes and Shortcomings of the Advanced Model}
\paragraph{Where the Advanced Model Succeeded:}
\begin{itemize}
    \item \textbf{Generalization to Rare Diagnoses:} The combination of balanced class weights and targeted dictionary augmentation enabled the model to correctly identify rare diagnostic classes. 
    \vspace{\baselineskip}

    \item \textbf{Disambiguating Procedural vs. Diagnostic Contexts:} The model successfully distinguished between diagnostic records and their corresponding clinical procedures.
    
    \item \textbf{Robustness to Lexical Noise:} Attention pooling allowed the classification head to bypass syntactic noise and focus on key clinical terms.
\end{itemize}

\paragraph{Where the Advanced Model didn't succeeded:}
\begin{itemize}
    \item \textbf{Ambiguity in Ultra-Short Literals:} Extremely brief literals consisting of a single abbreviation remained a major bottleneck. For example, the literal \texttt{ia} can represent \textit{Insuficiencia Aórtica} or \textit{Inteligencia Artificial} or \textit{Inseminación Artificial}. The lack of context in some sentences made the model defaulted to the maternally-biased class \texttt{O} or procedure class \texttt{0}, leading to systematic errors.
    
    \item \textbf{Granular boundary confusion:} The model occasionally struggled with the boundary between specialized endocrinology and general metabolic symptoms. Literals like \texttt{Deshidratación hipernatrémica} were misclassified as general symptoms instead of metabolic disorders .
\end{itemize}


% =========================================================================
% SECTION 4: Conclusions and Future Work
% =========================================================================
\section{Conclusions and Future Work}
In this academic project, our group successfully designed, optimized, and deployed an end-to-end clinical text classification pipeline for ICD-10 codification. By leveraging a specialized Spanish biomedical pre-trained backbone, designing a custom attention pooling mechanism, and implementing a strict stratified split-before-augment strategy, we transformed a difficult baseline into a highly robust clinical model capable of generalizing to rare diagnostic and procedural categories.

\subsection{Collaborative and Engineering Takeaways}
This project taught us that brute-force scaling of deep learning models is sub-optimal when working with real-world medical data and restricted academic hardware. The key engineering takeaways from our collaboration are:
\begin{itemize}
    \item Data quality and separation are essential: setting up a clean, leak-free validation split before running dictionary augmentations was the single most important factor in obtaining reliable evaluation metrics.
    
    \item Custom pooling and regularization, like Multi-Sample Dropout, provide substantial performance gains over simple pooling baselines, proving that custom network heads are vital when fine-tuning general encoders on domain-specific, short literals.
    
    \item Proactive hardware optimization is an essential skill in  AI engineering, allowing high-quality research to progress even under conditions like the ones we met in this project.
\end{itemize}

\subsection{Sociotechnical and Ethical Implications}
Deploying machine learning models in clinical environments carries significant ethical responsibility.
For instance, the bias and underrepresentations lead our training data to be heavily biased toward common hospital encounters. A model trained on this data might underperform on marginalized demographics or rare medical conditions, potentially leading to administrative misclassifications.

\end{document}
